\documentclass[11pt,a4paper]{article}
\usepackage[italian]{babel}
\usepackage[margin=2.5cm]{geometry}
\usepackage{parskip}
\usepackage{amsmath, amssymb, amsthm}
\usepackage{booktabs}
\usepackage{array}
\usepackage{xcolor}
\usepackage{needspace}
\usepackage{enumitem}
\usepackage{listings}
\usepackage{tikz}
\usepackage[most]{tcolorbox}
\usepackage[hidelinks]{hyperref}
\usetikzlibrary{decorations.pathreplacing, shapes.geometric, arrows.meta, positioning}

% Niente vedove/orfane: i paragrafi non lasciano righe isolate a fine/inizio pagina
\widowpenalty=10000
\clubpenalty=10000

% Elenchi più compatti
\setlist{itemsep=2pt, parsep=0pt, topsep=4pt, partopsep=0pt}

% --- Riquadri con bordo nero, senza numero (si spezzano tra due pagine) ---
% Uso: \begin{definizione} ... \end{definizione}
%      \begin{teorema}[Nome] ... \end{teorema}  (il nome è facoltativo)
\tcbset{nota/.style={enhanced jigsaw, breakable, colback=white,
  colframe=black, boxrule=0.5pt, sharp corners}}
\NewTColorBox{definizione}{}{nota,
  before upper={\textbf{Definizione.}\ }}
\NewTColorBox{teorema}{o}{nota,
  before upper={\textbf{Teorema\IfValueT{#1}{ (#1)}.}\ }}
\NewTColorBox{proposizione}{o}{nota,
  before upper={\textbf{Proposizione\IfValueT{#1}{ (#1)}.}\ }}
\NewTColorBox{proprieta}{o}{nota,
  before upper={\textbf{Proprietà\IfValueT{#1}{ (#1)}.}\ }}
\NewTColorBox{assioma}{o}{nota, boxrule=1pt,
  before upper={\textbf{Assioma\IfValueT{#1}{ (#1)}.}\ }}

% --- Ambienti semplici (senza riquadro) ---
% Il titolo sta in cima al blocco, anche se il contenuto inizia con un riquadro o
% con codice; e non resta mai da solo in fondo a una pagina.
% Uso: \begin{esempio}[Titolo facoltativo] ... \end{esempio}
\newcommand{\newrunin}[2]{%
  \NewDocumentEnvironment{#1}{o}{%
    \par\Needspace{4\baselineskip}\addvspace{\medskipamount}\noindent
    \textbf{#2}\IfValueT{##1}{ (##1)}\textbf{.}\ \ignorespaces}%
  {\par\addvspace{\medskipamount}}}
\newrunin{esempio}{Esempio}
\newrunin{esercizio}{Esercizio}
\newrunin{osservazione}{Osservazione}

% V e F nelle tavole di verità
\newcommand{\V}{\textsf{V}}
\newcommand{\F}{\textsf{F}}

% --- Codice C ---
% Uso: \begin{codice} ... \end{codice}      codice C numerato
%      \begin{comando} ... \end{comando}    comandi da terminale
%      \begin{risultato} ... \end{risultato} output di un programma
%      \lstinline|printf("ciao");|           codice dentro al testo
\lstdefinestyle{c}{
  language=C,
  basicstyle=\ttfamily\small,
  keywordstyle=\color{blue}\bfseries,
  commentstyle=\color{gray}\itshape,
  stringstyle=\color{red!70!black},
  numbers=left,
  numberstyle=\tiny\color{gray},
  numbersep=8pt,
  columns=flexible,
  keepspaces=true,
  breaklines=true,
  showstringspaces=false,
  tabsize=4,
  morekeywords={include, define},
  % lettere accentate nei commenti
  literate={à}{{\`a}}1 {è}{{\`e}}1 {é}{{\'e}}1 {ì}{{\`i}}1
           {ò}{{\`o}}1 {ù}{{\`u}}1 {È}{{\`E}}1 {À}{{\`A}}1
}
\lstset{style=c}
\newtcblisting{codice}{listing only, nota, left=6mm,
  listing options={style=c}}
\newtcblisting{comando}{listing only, nota,
  listing options={basicstyle=\ttfamily\small, numbers=none, language={},
    columns=flexible, keepspaces=true, breaklines=true}}
\newtcblisting{risultato}{listing only, enhanced jigsaw, breakable,
  colback=black!5, colframe=black, boxrule=0.5pt, sharp corners,
  listing options={basicstyle=\ttfamily\small, numbers=none, language={},
    columns=flexible, keepspaces=true, breaklines=true}}

% --- Diagrammi di flusso (simboli standard) ---
\tikzset{
  terminale/.style = {ellipse, draw, minimum width=2.5cm, minimum height=0.9cm},
  io/.style        = {trapezium, trapezium left angle=70, trapezium right angle=110,
                      draw, minimum width=2.5cm, minimum height=0.9cm},
  processo/.style  = {rectangle, draw, minimum width=2.5cm, minimum height=0.9cm},
  decisione/.style = {diamond, draw, aspect=2, minimum width=2.5cm},
  freccia/.style   = {thick, -{Stealth}}
}

\title{Limiti: monotone, algebra, potenze e polinomi}
\author{Marco Cerbella}
\date{Lezione 11 -- 29 ottobre 2026}

\begin{document}
\maketitle

\section{Limiti di funzioni monotone}

Ricordiamo: $f$ è \emph{monotona crescente} se $\forall x_1, x_2 \in D$, $x_1 < x_2 \Rightarrow f(x_1) \le f(x_2)$; \emph{monotona decrescente} se $x_1 < x_2 \Rightarrow f(x_1) \ge f(x_2)$.

\begin{proposizione}[Limite destro]
Sia $f$ una funzione definita e monotona in un intorno destro $I^+(c)$ del punto $c$ (dove $c$ è un numero reale oppure $-\infty$) escluso al più il punto $c$ stesso. Allora esiste, finito o infinito, il limite destro per $x \to c^+$ e precisamente si ha
\[
\lim_{x\to c^+} f(x) = \begin{cases} \inf\{f(x) : x \in I^+(c),\ x > c\} & \text{se } f \text{ è crescente} \\ \sup\{f(x) : x \in I^+(c),\ x > c\} & \text{se } f \text{ è decrescente.} \end{cases}
\]
\end{proposizione}

\begin{proposizione}[Limite sinistro]
Sia $f$ una funzione definita e monotona in un intorno sinistro $I^-(c)$ del punto $c$ (dove $c$ è un numero reale oppure $+\infty$) escluso al più il punto $c$ stesso. Allora esiste, finito o infinito, il limite sinistro per $x \to c^-$ e precisamente si ha
\[
\lim_{x\to c^-} f(x) = \begin{cases} \sup\{f(x) : x \in I^-(c),\ x < c\} & \text{se } f \text{ è crescente} \\ \inf\{f(x) : x \in I^-(c),\ x < c\} & \text{se } f \text{ è decrescente.} \end{cases}
\]
\end{proposizione}

Per una funzione crescente, il limite destro in $c$ è l'estremo inferiore dei valori a destra di $c$ (perché la funzione è limitata dal basso da $f(c)$ e $f$ diminuisce avvicinandosi a $c$); il limite sinistro è l'estremo superiore dei valori a sinistra. Se l'estremo inferiore è $-\infty$ il limite è $-\infty$, se l'estremo superiore è $+\infty$ il limite è $+\infty$.

\section{Continuità di funzioni monotone}

\begin{osservazione}
Una funzione monotona definita in un intorno di $x_0 \in \mathbb{R}$ può avere solo una discontinuità di prima specie.
\end{osservazione}

\begin{teorema}[Corollario]
Sia $f$ una funzione definita e monotona in un intorno $I(x_0)$ di un punto $x_0 \in \mathbb{R}$. Allora esistono finiti i limiti destro e sinistro per $x \to x_0$ e precisamente si ha:
\begin{enumerate}
    \item[(i)] se $f$ è crescente allora $\displaystyle\lim_{x\to x_0^-} f(x) \le f(x_0) \le \lim_{x\to x_0^+} f(x)$;
    \item[(ii)] se $f$ è decrescente allora $\displaystyle\lim_{x\to x_0^-} f(x) \ge f(x_0) \ge \lim_{x\to x_0^+} f(x)$.
\end{enumerate}
\end{teorema}

\section{Proprietà dei limiti}

\begin{teorema}[Unicità del limite]
Supponiamo che $f$ ammetta limite $\ell$ (finito o infinito) per $x$ tendente a $c$. Allora $f$ non ha altri limiti per $x$ tendente a $c$.
\end{teorema}

\begin{teorema}[Permanenza del segno]
Se per $x \to c$, $f(x) \to \ell > 0$, oppure $\ell = +\infty$, allora esiste un intorno $I(c)$ di $c$ tale che $f(x) > 0$ per $x \in I(c) \setminus \{c\}$. Un risultato analogo vale per il segno negativo.
\end{teorema}

\begin{teorema}[Inverso della permanenza del segno (1\textsuperscript{a} forma)]
Supponiamo che $f$ ammetta limite $\ell$ (finito o infinito) per $x$ tendente a $c$. Se esiste un intorno $I(c)$ di $c$ tale che $f(x) \ge 0$ in $I(c) \setminus \{c\}$ allora $\ell \ge 0$ oppure $\ell = +\infty$. Un risultato analogo vale per il segno negativo.
\end{teorema}

\section{Algebra dei limiti}

\begin{teorema}[Algebra dei limiti]
Se per $x \to x_0$, $f(x) \to \ell_1$ e $g(x) \to \ell_2$, $\ell_1, \ell_2 \in \mathbb{R}$, allora per $x \to x_0$ si ha:
\begin{itemize}
    \item $f(x) \pm g(x) \to \ell_1 \pm \ell_2$;
    \item $f(x)\,g(x) \to \ell_1 \ell_2$;
    \item $f(x)/g(x) \to \ell_1/\ell_2$ (purché $g(x), \ell_2 \neq 0$);
    \item $f(x)^{g(x)} \to \ell_1^{\ell_2}$ (purché $f(x), \ell_1 > 0$).
\end{itemize}
\end{teorema}

\begin{proprieta}[Aritmetizzazione parziale di $\infty$]
Per $a \in \mathbb{R}$:
\begin{itemize}
    \item $a \pm \infty = \pm\infty$; \quad $+\infty + \infty = +\infty$; \quad $-\infty - \infty = -\infty$;
    \item $a \cdot \infty = \infty$ ($a \neq 0$); \quad $\infty\cdot\infty = \infty$; \quad $\dfrac{a}{\infty} = 0$; \quad $\dfrac{\infty}{a} = \infty$ ($a \neq 0$).
\end{itemize}
Il segno si stabilisce con la regola dei segni.
\end{proprieta}

\begin{definizione}[Forme indeterminate]
Le forme indeterminate sono
\[
+\infty - \infty, \qquad 0 \cdot \infty, \qquad \frac{0}{0}, \qquad \frac{\infty}{\infty}, \qquad 1^{\pm\infty}, \qquad 0^0, \qquad (+\infty)^0 .
\]
\end{definizione}

\section{Limiti delle funzioni potenza}

Sia $f(x) = x^\alpha$.

Se $\alpha \in \mathbb{R}\setminus\mathbb{Q}$ allora $D(f) = (0, +\infty)$ e
\[
\lim_{x\to+\infty} x^\alpha = \begin{cases} +\infty & \alpha > 0 \\ 0 & \alpha < 0 \end{cases}
\qquad
\lim_{x\to 0^+} x^\alpha = \begin{cases} 0 & \alpha > 0 \\ +\infty & \alpha < 0. \end{cases}
\]

Se $f(x) = x^{m/n}$ con $n > 0$, $m \neq 0$, $m, n$ interi primi fra loro:
\begin{itemize}
    \item $n$ dispari e $m > 0$: $D(f) = \mathbb{R}$, $\displaystyle\lim_{x\to+\infty} x^{m/n} = +\infty$, $\displaystyle\lim_{x\to-\infty} x^{m/n} = \begin{cases} +\infty & m \text{ pari} \\ -\infty & m \text{ dispari} \end{cases}$ \quad (es.\ $\sqrt[3]{x^5}$);
    \item $n$ pari e $m > 0$: $D(f) = [0, +\infty)$, $\displaystyle\lim_{x\to+\infty} x^{m/n} = +\infty$ (es.\ $\sqrt{x}$);
    \item $n$ dispari e $m < 0$: $D(f) = \mathbb{R}\setminus\{0\}$, $\displaystyle\lim_{x\to\pm\infty} x^{m/n} = 0$, e per $x \to 0$: $\displaystyle\lim_{x\to 0^+} x^{m/n} = +\infty$, $\displaystyle\lim_{x\to 0^-} x^{m/n} = \begin{cases} +\infty & m \text{ pari} \\ -\infty & m \text{ dispari;} \end{cases}$
    \item $n$ pari e $m < 0$: $D(f) = \,]0, +\infty)$, $\displaystyle\lim_{x\to+\infty} x^{m/n} = 0$, $\displaystyle\lim_{x\to 0^+} x^{m/n} = +\infty$ (es.\ $\frac{1}{\sqrt x}$).
\end{itemize}

\section{Limiti di polinomi per $x \to \pm\infty$}

Per calcolare il limite di un polinomio per $x \to \pm\infty$ basta calcolare il limite del termine di grado massimo. Sia $P(x) = a_0 x^n + a_1 x^{n-1} + \dots + a_n$, $a_0 \neq 0$; si ha
\[
P(x) = a_0 x^n \left( 1 + \frac{a_1}{a_0 x} + \frac{a_2}{a_0 x^2} + \dots + \frac{a_n}{a_0 x^n} \right),
\qquad
\lim_{x\to\pm\infty} P(x) = \lim_{x\to\pm\infty} a_0 x^n ,
\]
infatti, tutti i termini dopo $1$, in parentesi graffa, tendono a zero per $x \to \pm\infty$ (e $1 + 0 + \dots + 0 = 1$).

\section{Esempi}

\begin{esempio}[Limite della tangente in $\frac{\pi}{2}$]
$\tan x = \dfrac{\sin x}{\cos x}$, $\sin\frac{\pi}{2} = 1$, $\cos\frac{\pi}{2} = 0$. Quindi $\displaystyle\lim_{x\to\frac{\pi}{2}} \tan x$ è del tipo $\frac{1}{0} = \infty$: bisogna studiare il segno.
\begin{itemize}
    \item $\cos x > 0$ se $x \in I^-(\frac{\pi}{2})$, $\cos x < 0$ se $x \in I^+(\frac{\pi}{2})$.
    \item Se $x \to \frac{\pi}{2}^-$: $\displaystyle\lim_{x\to\frac{\pi}{2}^-}\frac{\sin x}{\cos x} = +\infty$, perché $f(x) = \frac{\sin x}{\cos x} > 0$ per $x$ vicino a $\frac{\pi}{2}$ con $x < \frac{\pi}{2}$ (numeratore $>0$, denominatore $>0$).
    \item Se $x \to \frac{\pi}{2}^+$: $\displaystyle\lim_{x\to\frac{\pi}{2}^+}\frac{\sin x}{\cos x} = -\infty$, perché per $x$ vicino a $\frac{\pi}{2}$ con $x > \frac{\pi}{2}$ si ha numeratore $>0$ e denominatore $<0$.
\end{itemize}
Dunque $\displaystyle\nexists \lim_{x\to\frac{\pi}{2}}\tan x$, perché $\displaystyle\lim_{x\to\frac{\pi}{2}^-}\tan x \neq \lim_{x\to\frac{\pi}{2}^+}\tan x$.
\end{esempio}

\begin{esempio}[Polinomi e potenze]
\begin{itemize}
    \item $\displaystyle\lim_{x\to+\infty} (x^2 + x + 1) = \lim_{x\to+\infty} x^2\left(1 + \frac{1}{x} + \frac{1}{x^2}\right) = +\infty\cdot(1 + 0 + 0) = +\infty$, perché $\frac{1}{x} = x^{-1}$ e $\frac{1}{x^2} = x^{-2}$ sono potenze con esponente negativo, quindi tendono a $0$ (siamo nel caso $\frac{a}{\infty} = 0$) e $\lim_{x\to+\infty} x^2 = +\infty\cdot+\infty = +\infty$.
    \item $\displaystyle\lim_{x\to-\infty} (x^2 + x + 1) = \lim_{x\to-\infty} x^2\left(1 + \frac{1}{x} + \frac{1}{x^2}\right) = +\infty\cdot 1 = +\infty$.
    \item $\displaystyle\lim_{x\to 0} (x^2 + x + 1) = 1$, perché $x^2 + x + 1$ è una funzione continua in $\mathbb{R}$, definita in $x_0 = 0$, quindi il limite si ottiene sostituendo ad $x$ il valore $0$.
    \item $\displaystyle\lim_{x\to+\infty} (x + \sqrt{x}) = \lim_{x\to+\infty} x\left(1 + \frac{\sqrt{x}}{x}\right) = \lim_{x\to+\infty} x\left(1 + \frac{1}{\sqrt{x}}\right) = +\infty\cdot(1 + 0) = +\infty$, perché $\frac{\sqrt x}{x} = \frac{1}{\sqrt x} \to 0$ (siamo nella situazione $\frac{a}{\infty}$).
    \item $\displaystyle\lim_{x\to-\infty} (x + \sqrt{x})$ \emph{non ha senso}, perché $\sqrt{x}$ è definita solo per $x \ge 0$.
\end{itemize}
\end{esempio}

\end{document}
